\documentclass{amsart}
% For dealing with references we use the comment environment
\usepackage{verbatim}
\newenvironment{reference}{\comment}{\endcomment}
%\newenvironment{reference}{}{}
\newenvironment{slogan}{\comment}{\endcomment}
\newenvironment{history}{\comment}{\endcomment}

% For commutative diagrams we use Xy-pic
\usepackage[all]{xy}

% We use 2cell for 2-commutative diagrams.
\xyoption{2cell}
\UseAllTwocells

% We use multicol for the list of chapters between chapters
\usepackage{multicol}

% This is generally recommended for better output
\usepackage{lmodern}
\usepackage[T1]{fontenc}

% For cross-file-references

% Package for hypertext links:
\usepackage{hyperref}

% For any local file, say "hello.tex" you want to link to please

% Theorem environments.
%
\theoremstyle{plain}
\newtheorem{theorem}[subsection]{Theorem}
\newtheorem{proposition}[subsection]{Proposition}
\newtheorem{lemma}[subsection]{Lemma}

\theoremstyle{definition}
\newtheorem{definition}[subsection]{Definition}
\newtheorem{example}[subsection]{Example}
\newtheorem{exercise}[subsection]{Exercise}
\newtheorem{situation}[subsection]{Situation}

\theoremstyle{remark}
\newtheorem{remark}[subsection]{Remark}
\newtheorem{remarks}[subsection]{Remarks}

\numberwithin{equation}{subsection}

% Macros
%
\def\lim{\mathop{\mathrm{lim}}\nolimits}
\def\colim{\mathop{\mathrm{colim}}\nolimits}
\def\Spec{\mathop{\mathrm{Spec}}}
\def\Hom{\mathop{\mathrm{Hom}}\nolimits}
\def\Ext{\mathop{\mathrm{Ext}}\nolimits}
\def\SheafHom{\mathop{\mathcal{H}\!\mathit{om}}\nolimits}
\def\SheafExt{\mathop{\mathcal{E}\!\mathit{xt}}\nolimits}
\def\Sch{\mathit{Sch}}
\def\Mor{\mathop{\mathrm{Mor}}\nolimits}
\def\Ob{\mathop{\mathrm{Ob}}\nolimits}
\def\Sh{\mathop{\mathit{Sh}}\nolimits}
\def\NL{\mathop{N\!L}\nolimits}
\def\CH{\mathop{\mathrm{CH}}\nolimits}
\def\proetale{{pro\text{-}\acute{e}tale}}
\def\etale{{\acute{e}tale}}
\def\QCoh{\mathit{QCoh}}
\def\Ker{\mathop{\mathrm{Ker}}}
\def\Im{\mathop{\mathrm{Im}}}
\def\Coker{\mathop{\mathrm{Coker}}}
\def\Coim{\mathop{\mathrm{Coim}}}

% Boxtimes
%
\DeclareMathSymbol{\boxtimes}{\mathbin}{AMSa}{"02}

%
% Macros for moduli stacks/spaces
%
\def\QCohstack{\mathcal{QC}\!\mathit{oh}}
\def\Cohstack{\mathcal{C}\!\mathit{oh}}
\def\Spacesstack{\mathcal{S}\!\mathit{paces}}
\def\Quotfunctor{\mathrm{Quot}}
\def\Hilbfunctor{\mathrm{Hilb}}
\def\Curvesstack{\mathcal{C}\!\mathit{urves}}
\def\Polarizedstack{\mathcal{P}\!\mathit{olarized}}
\def\Complexesstack{\mathcal{C}\!\mathit{omplexes}}
% \Pic is the operator that assigns to X its picard group, usage \Pic(X)
% \Picardstack_{X/B} denotes the Picard stack of X over B
% \Picardfunctor_{X/B} denotes the Picard functor of X over B
\def\Pic{\mathop{\mathrm{Pic}}\nolimits}
\def\Picardstack{\mathcal{P}\!\mathit{ic}}
\def\Picardfunctor{\mathrm{Pic}}
\def\Deformationcategory{\mathcal{D}\!\mathit{ef}}

\setlength{\emergencystretch}{3em}
\begin{document}
\title[Stacks Project, Tag 09NA]{583 --- Line bundles extend across discrete retrocompact complements with Noetherian local rings of dimension at most one}
\maketitle
\noindent Copyright (C) 2005--2025 Johan de Jong. Stacks Project authors. Source: \href{https://stacks.math.columbia.edu/tag/09NA}{Tag 09NA}.

\noindent Editorial difficulty note (not author text). Difficulty H1: The proof constructs local equations that isolate each missing point even when the ambient scheme is not globally Noetherian. It glues these equations to a line bundle realizing the open complement and uses triviality of Picard groups on localized Artinian rings to extend arbitrary line bundles from that open..

\noindent Editorial provenance note (not author text). History: excerpt from revision \texttt{a0444\allowbreak 6e57e\allowbreak c1fbc\allowbreak 252a8\allowbreak 71afc\allowbreak ec775\allowbreak 2fb28\allowbreak 07b14}, varieties.tex, lines 7796--7853. Reference presentation was adapted; original-author context and core support blocks were retained or added. The mathematical wording is unchanged. Licensed under GNU FDL 1.2 or later; no invariant sections or cover texts. See ../licenses/COPYING. Context blocks reproduce author definitions and supporting results; they are not additional counted problems.

\begin{lemma}
\label{lemma-complement-codim-1-closed-points}
Let $X$ be a scheme. Let $U \subset X$ be an open. Assume
\begin{enumerate}
\item $U$ is a retrocompact open of $X$,
\item $X \setminus U$ is discrete, and
\item for $x \in X \setminus U$ the local ring
$\mathcal{O}_{X, x}$ is Noetherian of dimension $\leq 1$.
\end{enumerate}
Then (1) there exists an invertible $\mathcal{O}_X$-module $\mathcal{L}$
and a section $s$ such that $U = X_s$ and (2) the map
$\Pic(X) \to \Pic(U)$ is surjective.
\end{lemma}

\begin{proof}
Let $X \setminus U = \{x_i; i \in I\}$.
Choose affine opens $U_i \subset X$ with $x_i \in U_i$ and
$x_j \not \in U_i$ for $j \not = i$. This is possible by condition (2).
Say $U_i = \Spec(A_i)$. Let $\mathfrak m_i \subset A_i$ be the maximal
ideal corresponding to $x_i$. By our assumption on the local rings
there are only a finite number of prime ideals
$\mathfrak q \subset \mathfrak m_i$,
$\mathfrak q \not = \mathfrak m_i$ (see
Algebra, Lemma \href{https://stacks.math.columbia.edu/tag/00FR}{lemma Noetherian irreducible components (Tag 00FR)}).
Thus by prime avoidance (Algebra, Lemma
\href{https://stacks.math.columbia.edu/tag/00DS}{lemma silly (Tag 00DS)}) we can find $f_i \in \mathfrak m_i$
not contained in any of those primes. Then
$V(f_i) = \{\mathfrak m_i\} \amalg Z_i$ for some closed subset
$Z_i \subset U_i$ because $Z_i$ is a retrocompact open subset of
$V(f_i)$ closed under specialization, see
Algebra, Lemma \href{https://stacks.math.columbia.edu/tag/00I0}{lemma constructible stable specialization closed (Tag 00I0)}.
After shrinking $U_i$ we may assume $V(f_i) = \{x_i\}$. Then
$$
\mathcal{U} : X = U \cup \bigcup U_i
$$
is an open covering of $X$. Consider the $2$-cocycle with values
in $\mathcal{O}_X^*$ given by $f_i$ on $U \cap U_i$ and by
$f_i/f_j$ on $U_i \cap U_j$. This defines a line bundle
$\mathcal{L}$ such that the section $s$ defined by $1$ on $U$
and $f_i$ on $U_i$ is as in the statement of the lemma.

\medskip\noindent
Let $\mathcal{N}$ be an invertible $\mathcal{O}_U$-module.
Let $N_i$ be the invertible $(A_i)_{f_i}$ module such that
$\mathcal{N}|_{U \cap U_i}$ is equal to $\tilde N_i$.
Observe that $(A_{\mathfrak m_i})_{f_i}$ is an Artinian ring
(as a dimension zero Noetherian ring, see
Algebra, Lemma \href{https://stacks.math.columbia.edu/tag/00KH}{lemma Noetherian dimension 0 (Tag 00KH)}).
Thus it is a product of local rings
(Algebra, Lemma \href{https://stacks.math.columbia.edu/tag/00JB}{lemma artinian finite length (Tag 00JB)}) and
hence has trivial Picard group. Thus, after shrinking $U_i$
(i.e., after replacing $A_i$ by $(A_i)_g$ for some $g \in A_i$,
$g \not \in \mathfrak m_i$)
we can assume that $N_i = (A_i)_{f_i}$, i.e., that
$\mathcal{N}|_{U \cap U_i}$ is trivial. In this case it is
clear how to extend $\mathcal{N}$ to an invertible sheaf over $X$
(by extending it by a trivial invertible module over each $U_i$).
\end{proof}
\end{document}
